% Part: normal-modal-logic
% Chapter: tableaux
% Section: rules-for-K

\documentclass[../../../include/open-logic-section]{subfiles}

\begin{document}

\olfileid{nml}{tab}{rul}

\olsection{\Ax{K}-এর বিধি}

সাধারণ বচনসংযোজকগুলির বিধি, পূর্বসূচি যোগ করা ছাড়া, সাধারণ চিহ্নিত
বচনমূলক !!{tableau}s-এর বিধিরই অনুরূপ। প্রতিটি ক্ষেত্রে চিহ্নিত
!!{formula} $\sFmla{S}{!A}[\sigma]$-তে বিধি প্রয়োগ করলে যে নতুন
!!{formula}s পাওয়া যায়, তাদেরও পূর্বসূচি~$\sigma$। এটি স্বতঃবোধ্য:
যেমন, $\sigma$-র নামাঙ্কিত জগতে $!A \land !B$ সত্য হলে, $\sigma$-তেই
$!A$ ও $!B$ সত্য (অন্য কোনো জগৎ সম্পর্কে এতে কিছু বলা হয় না)।
বচনমূলক বিধিগুলি \olref{tab:prop-rules}-এ একত্র দেওয়া হয়েছে।

\begin{table}
  \[\def\arraystretch{3}\begin{array}{|c|c|}
    \hline
    \AxiomC{\sFmla{\True}{\lnot !A}[\sigma]}
    \RightLabel{\TRule{\True}{\lnot}}
    \UnaryInfC{\sFmla{\False}{!A}[\sigma]}
    \DisplayProof
    &
    \AxiomC{\sFmla{\False}{\lnot !A}[\sigma]}
    \RightLabel{\TRule{\False}{\lnot}}
    \UnaryInfC{\sFmla{\True}{!A}[\sigma]}
    \DisplayProof
    \\[1ex]
    \hline
    \AxiomC{\sFmla{\True}{!A \land !B}[\sigma]}
    \RightLabel{\TRule{\True}{\land}}
    \UnaryInfC{\sFmla{\True}{!A}[\sigma]}
    \noLine
    \UnaryInfC{\sFmla{\True}{!B}[\sigma]}
    \DisplayProof
    &
    \AxiomC{\sFmla{\False}{!A \land !B}[\sigma]}
    \RightLabel{\TRule{\False}{\land}}
    \UnaryInfC{$\sFmla{\False}{!A}[\sigma] \quad \mid \quad
      \sFmla{\False}{!B}[\sigma]$}
    \DisplayProof
    \\[2ex]
    \hline
    \AxiomC{\sFmla{\True}{!A \lor !B}[\sigma]}
    \RightLabel{\TRule{\True}{\lor}}
    \UnaryInfC{$\sFmla{\True}{!A}[\sigma] \quad \mid \quad
      \sFmla{\True}{!B}[\sigma]$}
    \DisplayProof
    &
    \AxiomC{\sFmla{\False}{!A \lor !B}[\sigma]}
    \RightLabel{\TRule{\False}{\lor}}
    \UnaryInfC{\sFmla{\False}{!A}[\sigma]}
    \noLine
    \UnaryInfC{\sFmla{\False}{!B}[\sigma]}
    \DisplayProof
    \\[2ex]
    \hline
    \AxiomC{\sFmla{\True}{!A \lif !B}[\sigma]}
    \RightLabel{\TRule{\True}{\lif}}
    \UnaryInfC{$\sFmla{\False}{!A}[\sigma] \quad \mid
      \quad \sFmla{\True}{!B}[\sigma]$}
    \DisplayProof
    &
    \AxiomC{\sFmla{\False}{!A \lif !B}[\sigma]}
    \RightLabel{\TRule{\False}{\lif}}
    \UnaryInfC{\sFmla{\True}{!A}[\sigma]}
    \noLine
    \UnaryInfC{\sFmla{\False}{!B}[\sigma]}
    \DisplayProof
    \\[2ex]
    \hline
  \end{array}\]
  \caption{বচনসংযোজকগুলির জন্য পূর্বসূচিযুক্ত !!{tableau}-র বিধি}
  \ollabel{tab:prop-rules}
\end{table}

সাধারণ !!{tableau}s-এর মতোই এখানে বদ্ধতার শর্ত, তবে শুধু
!!{formula}s নয়, পূর্বসূচিও একই হতে হবে। তাই কোনো শাখায়
\[
\sFmla{\True}{!A}[\sigma] \quad\text{ও}\quad \sFmla{\False}{!A}[\sigma]
\]
উভয়ই থাকলে, কোনো পূর্বসূচি $\sigma$ ও !!{formula}~$!A$-এর জন্য,
শাখাটি বদ্ধ।

অনুমিতি স্থাপনের বিধিও সাধারণ !!{tableau}s-এর মতোই, শুধু অনুমিতির
পূর্বসূচি হিসেবে আমরা সর্বদা~$1$ ব্যবহার করি। (সব অনুমিতিতে একই
পূর্বসূচি ব্যবহার করলেই চলে; সেটি কোন পূর্বসূচি, তা গুরুত্বপূর্ণ নয়।)
সুতরাং, যেমন,
\[
!B_1, \dots, !B_n \Proves !A
\]
তখনই এবং কেবল তখনই, যখন নিচের অনুমিতিগুলির একটি বদ্ধ ট্যাবলো আছে:
\[
\sFmla{\True}{!B_1}[1], \dots, \sFmla{\True}{!B_n}[1],
\sFmla{\False}{!A}[1].
\]

মোডাল অপারেটর\iftag{prvBox}{\iftag{prvDiamond}{গুলি~$\Box$
    ও}{~$\Box$}}{}\iftag{prvDiamond}{~$\Diamond$}{}-এর ক্ষেত্রে
পূর্বসূচি~$\sigma$-যুক্ত !!a{formula}-তে বিধি প্রয়োগের সিদ্ধান্তের
পূর্বসূচি হয় $\sigma.n$। তবে কোন $n$ ব্যবহার করা যাবে, তা চিহ্নটি
$\True$ না~$\False$, তার উপর নির্ভর করে।

\iftag{prvBox}{$\TRule{\True}{\Box}$ বিধি
  $\sFmla{\True}{\Box !A}[\sigma]$-যুক্ত শাখায়
  $\sFmla{\True}{!A}[\sigma.n]$ যোগ করে।
  \iftag{prvDiamond}{একইভাবে,}{}}{}
\iftag{prvDiamond}{$\TRule{\False}{\Diamond}$ বিধি
  $\sFmla{\False}{\Diamond !A}[\sigma]$-যুক্ত শাখায়
  $\sFmla{\False}{!A}[\sigma.n]$ যোগ করে।}{}
এই বিধি\iftag{notprvBox,notprvDiamond}{}{গুলি} কেবল তখনই প্রয়োগ
করা যায়, যখন পূর্বসূচি~$\sigma.n$ সংশ্লিষ্ট শাখায় \emph{আগেই}
আছে। এমন পূর্বসূচিকে সেই শাখায় ‘ব্যবহৃত’ বলব।

\iftag{prvBox}{$\TRule{\False}{\Box}$ বিধি
  $\sFmla{\False}{\Box !A}[\sigma]$-যুক্ত শাখায়
  $\sFmla{\False}{!A}[\sigma.n]$ যোগ করে।
  \iftag{prvDiamond}{একইভাবে,}{}}{}
\iftag{prvDiamond}{$\TRule{\True}{\Diamond}$ বিধি
  $\sFmla{\True}{\Diamond !A}[\sigma]$-যুক্ত শাখায়
  $\sFmla{\True}{!A}[\sigma.n]$ যোগ করে।}{}
কিন্তু এই বিধি\iftag{notprvBox,notprvDiamond}{}{গুলি} কেবল তখনই
প্রয়োগ করা যায়, যখন পূর্বসূচি~$\sigma.n$ সংশ্লিষ্ট শাখায়
\emph{আগে ছিল না}। এমন পূর্বসূচিকে শাখায় ‘নতুন’ বলি।

বিধিগুলি \olref{tab:rules-K}-এ দেওয়া হয়েছে।

\begin{table}
  \begin{center}
    \def\arraystretch{3}\def\fCenter{}
    \begin{tabular}{|c|c|}
    \hline
    \iftag{prvBox}{
      \AxiomC{\sFmla{\True}{\Box !A}[\sigma]}
      \RightLabel{\TRule{\True}{\Box}}
      \UnaryInfC{\sFmla{\True}{!A}[\sigma.n]}
      \DisplayProof
      &
      \AxiomC{\sFmla{\False}{\Box !A}[\sigma]}
      \RightLabel{\TRule{\False}{\Box}}
      \UnaryInfC{\sFmla{\False}{!A}[\sigma.n]}
      \DisplayProof\\
      $\sigma.n$ ব্যবহৃত & $\sigma.n$ নতুন
      \\[1ex]
      \hline}{}
    \iftag{prvDiamond}{
      \AxiomC{\sFmla{\True}{\Diamond !A}[\sigma]}
      \RightLabel{\TRule{\True}{\Diamond}}
      \UnaryInfC{\sFmla{\True}{!A}[\sigma.n]}
      \DisplayProof
      &
      \AxiomC{\sFmla{\False}{\Diamond !A}[\sigma]}
      \RightLabel{\TRule{\False}{\Diamond}}
      \UnaryInfC{\sFmla{\False}{!A}[\sigma.n]}
      \DisplayProof\\
      $\sigma.n$ নতুন & $\sigma.n$ ব্যবহৃত
      \\[1ex]
      \hline}{}
    \end{tabular}
  \end{center}
  \caption{\Ax{K}-এর মোডাল বিধি।}
  \ollabel{tab:rules-K}
\end{table}

\iftag{prvBox}{\TRule{\True}{\Box}}{\TRule{\False}{\Diamond}}-এর
পূর্বসূচি ব্যবহৃত হতে হবে—এই বিধিনিষেধ জরুরি। তা না হলে নিচেরটিকে
বদ্ধ !!{tableau} হিসেবে গণ্য করতে হতো:
\iftag{notprvBox,notprvDiamond}{%
  \iftag{prvBox}{%
  \begin{oltableau}
    [\pFmla{\True}{\Box \formula{A}}{1}, just = \TAss
      [\pFmla{\False}{\lnot\Box\lnot \formula{A}}{1}, just = \TAss
        [\pFmla{\True}{\formula{A}}{1.1}, just = {\TRule{\True}{\Box}[1]}
          [\pFmla{\True}{\Box\lnot\formula{A}}{1},
            just ={\TRule{\False}{\lnot}[2]}
            [\pFmla{\True}{\lnot\formula{A}}{1.1},
              just = {\TRule{\True}{\Box}[4]}
              [\pFmla{\False}{\formula{A}}{1.1},
                  just ={\TRule{\True}{\lnot}[5]}, close]
            ]
          ]
        ]
      ]
    ]
  \end{oltableau}
  }{
  \begin{oltableau}
    [\pFmla{\True}{\lnot\Diamond\lnot \formula{A}}{1}, just = \TAss
      [\pFmla{\False}{\Diamond\formula{A}}{1}, just = \TAss
        [\pFmla{\False}{\formula{A}}{1.1}, just = {\TRule{\False}{\Diamond}[2]}
          [\pFmla{\False}{\Diamond\lnot\formula{A}}{1},
            just ={\TRule{\True}{\lnot}[1]}
            [\pFmla{\False}{\lnot\formula{A}}{1.1},
              just = {\TRule{\False}{\Diamond}[4]}
              [\pFmla{\True}{\formula{A}}{1.1},
                  just ={\TRule{\True}{\lnot}[5]}, close]
            ]
          ]
        ]
      ]
    ]
  \end{oltableau}
}}{
  \begin{oltableau}
    [\pFmla{\True}{\Box \formula{A}}{1}, just = \TAss
      [\pFmla{\False}{\Diamond \formula{A}}{1}, just = \TAss
        [\pFmla{\True}{\formula{A}}{1.1}, just = {\TRule{\True}{\Box}[1]}
          [\pFmla{\False}{\formula{A}}{1.1},
            just = {\TRule{\False}{\Diamond}[2]}, close]
        ]
      ]
    ]
\end{oltableau}}

কিন্তু $\Box \formula{A} \Entails/ \Diamond \formula{A}$; তাই আমাদের
প্রমাণ-পদ্ধতি অবিশুদ্ধ হয়ে পড়ত। একইভাবে,
$\Diamond \formula{A} \Entails/ \Box \formula{A}$; কিন্তু
\iftag{prvBox}{\TRule{\False}{\Box}}{\TRule{\True}{\Diamond}}-এর
পূর্বসূচি নতুন হতে হবে—এই বিধিনিষেধ বাদ দিলে নিচেরটিও একটি বদ্ধ
ট্যাবলো হতো: \iftag{notprvBox,notprvDiamond}{%
  \iftag{prvBox}{%
  \begin{oltableau}
    [\pFmla{\True}{\lnot\Box\lnot \formula{A}}{1}, just = \TAss
      [\pFmla{\False}{\Box \formula{A}}{1}, just = \TAss
        % BN-SRC-371: False box on line 2, not true box, supplies this conclusion.
        [\pFmla{\False}{\formula{A}}{1.1}, just = {\TRule{\False}{\Box}[2]}
          [\pFmla{\False}{\Box\lnot\formula{A}}{1},
            just ={\TRule{\True}{\lnot}[1]}
            [\pFmla{\False}{\lnot\formula{A}}{1.1},
              just = {\TRule{\False}{\Box}[4]}
              [\pFmla{\True}{\formula{A}}{1.1},
                  just ={\TRule{\False}{\lnot}[5]}, close]
            ]
          ]
        ]
      ]
    ]
  \end{oltableau}
  }{
  \begin{oltableau}
    [\pFmla{\True}{\Diamond \formula{A}}{1}, just = \TAss
      [\pFmla{\False}{\lnot\Diamond\lnot\formula{A}}{1}, just = \TAss
        [\pFmla{\True}{\formula{A}}{1.1}, just = {\TRule{\True}{\Diamond}[1]}
          [\pFmla{\True}{\Diamond\lnot\formula{A}}{1},
            just ={\TRule{\False}{\lnot}[2]}
            [\pFmla{\True}{\lnot\formula{A}}{1.1},
              just = {\TRule{\True}{\Diamond}[4]}
              [\pFmla{\False}{\formula{A}}{1.1},
                  just ={\TRule{\True}{\lnot}[5]}, close]
            ]
          ]
        ]
      ]
    ]
  \end{oltableau}
}}{
  \begin{oltableau}
    [\pFmla{\True}{\Diamond \formula{A}}{1}, just = \TAss,name=one
      [\pFmla{\False}{\Box \formula{A}}{1}, just = \TAss, name=two
        [\pFmla{\True}{\formula{A}}{1.1}, just = {\TRule{\True}{\Diamond}[1]}
          [\pFmla{\False}{\formula{A}}{1.1},
            just = {\TRule{\False}{\Box}[2]}, close]
        ]
      ]
    ]
  \end{oltableau}}

\end{document}
